\begin{frame} \frametitle{Navigation}
\vspace{-10pt}
\textit{\textbf{Navigation}} implies the capabilities of:

{\footnotesize $\bullet$ accurate determination of the vehicle position and velocity with respect to a known reference point}

{\footnotesize $\bullet$ planning and the execution of the movements between locations}
\vspace{5pt}
\begin{columns}
\column{.4\textwidth}
\centering
\includegraphics[width=0.8\linewidth]{fig/map.pdf} \\
{\scriptsize John Harrison, $18^{th}$ century \\
\textit{longitude problem} }


\column{.2\textwidth}
\centering
\includegraphics[height=5em]{fig/phone-navigation.pdf} \\
{\scriptsize GPS, USDOD, 1974-1994. }

\column{.3\textwidth}
\centering
\includegraphics[height=5em]{fig/rov-navigation.pdf} \\
{\scriptsize  Underwater navigation strategies }
\end{columns} 
\vspace{5pt}
\hspace{0.5cm} \pro Measured depth is fairly accurate $\rightarrow$ localisation mainly 2d

\hspace{0.5cm} \contra Obtaining absolute position is more complex underwater

\hspace{0.5cm} \contra Cameras and sonar are of a limited use

\hspace{0.5cm} \contra Dead-reckoning sensors are expensive and prone to drifting

{\centering 
$\Longrightarrow$ use all available inertial information to navigate
}
\end{frame}